\subsection{File System Data Structures}

\begin{lstlisting}[
    style=myagc,
    caption={File system description: \file{include/fs/fs\_desc.h}},
    linebackgroundcolor={\lines{18}{19} \lines{4}{5} \line{8} \line{10}},
    label={fs:desc}
]
template <...>
struct fs_desc_t {
  using addr_t = /* type of address supported by device */;
    // for flash, it is uint32_t since writes must be word aligned
    // for fram, it can be uint16_t to save space 

  size_t blk_sz; /* size of a block */
    // for flash, it is 1024, 
    //   since that's the smallest unit we can effectively write
    // for FRAM, it is 128

  size_t max_num_blks; /* man number of blocks per file */
  size_t n_inodes; /* number of files */

  size_t n_open_files; /* number of open files */


  using fs_xfer_t = size_t (*)(addr_t addr, uint8_t *buf, size_t buf_sz);

  // methods to transfer data to and from the target device
  fs_xfer_t write, write_sync;
  fs_xfer_t read, read_sync;
};
\end{lstlisting}


\begin{lstlisting}[
    style=myagc,
    caption={Filesystem backend: \file{include/fs/fs\_bck.h}},
    linebackgroundcolor={},
    label={fs:fs}
]
template <...>
class fs_t
{
  struct hd_t {
    inode_t tbl[N_INODES];
    uint32_t magic;
    bitset<N_BLKS> free_set;
  };

  inline static hd_t hd;

public:
  // return the corresponding inode, and whether the file was just opened
  static inline pair<const inode_t *, bool> open(fn_t fn, size_t sz,
                                                 uint32_t flags)
  {
    assert(fn >= 0 && fn < desc.n_inodes, TERM);

    auto &inode = hd.tbl[fn];
    const bool in_use = inode.in_use();
    const bool to_create = flags & O_CREATE;

    if (in_use || !to_create) {
      assert(in_use, TERM, "file does not exist\r\n");
      if (sz != 0)
        debug<TRACE>("passing size (%d) to open existing file %d\r\n", sz, fn);
      return {&inode, false};
    }

    assert(to_create, TERM, "not creating non-existent file %d\r\n", fn);
    inode.set_use();

    uint8_t nblks = roundup(max(sz, BLK_SZ), BLK_SZ) / BLK_SZ;
    assert(nblks <= desc.max_num_blks, TERM);
    const_cast<uint8_t &>(inode.nblks) = nblks;

    auto success = alloc_blks(inode);
    assert(success, TERM, "block allocation failed for file %d\r\n", fn);

    if constexpr (desc.is_ram)
      inode.set_ft(flags & O_CHAR_FILE ? CHAR : BIN);
    else
      inode.set_ft(BIN);

    if (flags & O_PERM)
      inode.set_perm();

    inode.end = 0;

    // partial update, only modifies the inode entry on the device
    update_hd(inode); 

    return {&inode, true};
  }

private:
  // return the number of bytes written/read
  template <auto func>
  static inline size_t xfer(const inode_t *inode, uint8_t *buf,
                            const size_t buf_sz, size_t offset = 0)
  {
    assert(inode->max_sz() >= offset + buf_sz, TERM);
    const auto nblks = inode->nblks;

    uint8_t fst_blk = offset / BLK_SZ;
    size_t blk_offset = offset % BLK_SZ;
    size_t fst_blk_sz = min(buf_sz, BLK_SZ - blk_offset);

    uint8_t i = fst_blk;

    size_t sent_sz;

    auto rem = buf_sz;
    sent_sz = func(paddr(inode->blks[i++]) + blk_offset, buf, fst_blk_sz);
    if (sent_sz < fst_blk_sz)
      return sent_sz;

    buf += fst_blk_sz;
    rem -= fst_blk_sz;

    while (rem > 0 && i < nblks) {
      auto blk_sz = min(rem, BLK_SZ);

      sent_sz = func(paddr(inode->blks[i++]), buf, blk_sz);
      if (sent_sz < blk_sz)
        return buf_sz - (rem - sent_sz);

      buf += blk_sz;
      rem -= blk_sz;
    }

    return buf_sz;
  }

public:
  template <bool sync = ASYNC>
  static size_t load(const inode_t *inode, void *buf, size_t buf_sz,
                     size_t off = 0)
  {
    constexpr auto read = sync ? desc.read_sync : desc.read;
    return xfer<read>(inode, (uint8_t *)buf, buf_sz, off);
  }

  template <bool sync = ASYNC>
  static size_t store(const inode_t *inode, const void *buf, size_t buf_sz,
                      size_t off = 0)
  {
    constexpr auto write = sync ? desc.write_sync : desc.write;
    auto nbytes = xfer<write>(inode, (uint8_t *)buf, buf_sz, off);

    bool prev_invalid = inode->end == 0;
    inode->end = max(inode->end, off + nbytes);

    return nbytes;
  }
};
\end{lstlisting}


\begin{lstlisting}[
    style=myagc,
    caption={Input file stream: \file{include/fs/ifstream.h}},
    linebackgroundcolor={},
    label={fs:ifstream}
]
template <auto desc>
class ifstream
{
  using file_t = fs_impl_t<desc>::file_t;
  file_t file;

  static constexpr size_t buf_len = 16;
  char buf[buf_len];
  size_t buf_pos = buf_len;

  size_t remaining;
  size_t idx;
  size_t offset;

  inline void fetch()
  {
    size_t n_chars = min(remaining, buf_len);
    file.read((uint8_t *)buf, n_chars, offset);

    remaining -= n_chars;
    offset += n_chars;
    buf_pos = 0;
  }

public:
  ifstream(fn_t fn, size_t sz, size_t off = 0)
      : file{fn, O_READ}, remaining{sz}, idx{sz}, offset{off}
  {
  }

  char operator*()
  {
    if (idx == 0)
      return '\0';

    if (buf_pos == buf_len)
      fetch();

    return buf[buf_pos];
  }

  ifstream &operator++()
  {
    if (buf_pos == buf_len)
      fetch();

    idx--;
    buf_pos++;
    return *this;
  }
\end{lstlisting}


\begin{lstlisting}[
    style=myagc,
    caption={Output file stream: \file{include/fs/ofstream.h}},
    linebackgroundcolor={},
    label={fs:ofstream}
]
template <auto desc>
class ofstream
{
  using file_t = fs_impl_t<desc>::file_t;
  file_t file;

  static constexpr size_t buf_len = 16;
  char buf[buf_len];
  size_t buf_pos = 0;

  size_t offset;

  void flush()
  {
    if (buf_pos != 0) {
      file.write(buf, buf_pos, offset);
      offset += buf_pos;
      buf_pos = 0;
    }
  }

public:
  ofstream(fn_t fn, uint32_t flags = 0)
      : file{fn, flags | O_WRITE | O_CHAR_FILE}, offset{0}
  {
  }

  void seek(size_t off) { offset = off == EOF ? file.fend() : off; }

  void putc(char c)
  {
    if (buf_pos == buf_len)
      flush();
    buf[buf_pos++] = c;
  }

  template <typename... Args>
  void printf(Args... args)
  {
    do_printf(*this, std::forward<Args>(args)...);
  }
};
\end{lstlisting}


\begin{lstlisting}[
    style=myagc,
    caption={File handler, file opening and closing sequence},
    linebackgroundcolor={}
]
class file_t {
  fd_t *fd = nullptr;
  /* memory mapping information */

public:

  file_t(fn_t fn, size_t sz, uint32_t flags) {
    fd = find_fd(fn); // first find available file descriptor

    // then open the file, getting the inode
    auto [inode, _new_file] = fs.open(fn, sz, flags);

    // if the file is already open
    if (fd->fn == fn) {
      assert(!_new_file && fd->inode == inode, H_RESET);

      fd->r_cnt++;
      fd->w_cnt += w_en;
      return;
    }

    // otherwise set up the file descriptor
    fd->fn = fn;
    fd->inode = inode;

    fd->mu = mmap_unit_t{};
    fd->r_cnt = 1;
    fd->w_cnt = w_en;
  }

  ~file_t() {
    if (fd != nullptr) // if it's not alrady closed
      close();
  }

  void close() {
    unmap();

    fd->r_cnt--;
    fd->w_cnt -= w_en;

    fs.close(fd->inode);

    // last reference
    if (fd->r_cnt == 0) {
      fd->fn = null_fn;
      fd->inode = nullptr;
      fd->mu = mmap_unit_t{};
    }

    fd = nullptr;
  }
};

\end{lstlisting}
